% dossier.tex — GENERATED by build.py; edit printed.yaml or intro.texfrag, then `make`.
% Public version.
\documentclass[11pt]{article}
\usepackage[letterpaper,margin=1in]{geometry}
\usepackage{fontspec}
\IfFileExists{libertinus.sty}{\usepackage{libertinus}}{\setmainfont{DejaVu Serif}}
\IfFileExists{danish.ldf}{\usepackage[english,danish]{babel}}{\usepackage[english]{babel}}
\usepackage{microtype,xcolor,booktabs,longtable,array}
\usepackage[hidelinks]{hyperref}
\setcounter{secnumdepth}{0}
\setcounter{tocdepth}{2}
\emergencystretch=5em
\tolerance=1500
\newcommand{\dcite}[1]{{\nopagebreak\raggedleft\small #1\par}\medskip}
\newcommand{\tlg}[1]{{\setlength{\fboxsep}{0.8pt}\colorbox{yellow}{#1}}}
\newcommand{\olg}[1]{{\setlength{\fboxsep}{0.8pt}\colorbox{cyan!25}{#1}}}
\title{\emph{Entydig} hos Niels Bohr\\[0.5ex]\large Kildedossier}
\author{Compiled from sks-search and bibliotek-search}
\date{26 September 2026}
\begin{document}
\maketitle
\begin{otherlanguage}{english}

\section{About this dossier}
\noindent Danish \tlg{entydig} is marked in yellow throughout; in the printed
originals, the corresponding \olg{eindeutig}, \olg{unique} or \olg{unambiguous}
is marked in blue. Bohr's writings are in copyright. \emph{Filosofiske Skrifter} I--IV appear in this collection with the rights-holders' permission; the other Bohr passages here are brief quotations for the purposes of study, and the aligned originals are cut to the sentences that bear on the word. Its four parts are listed in the contents.
\end{otherlanguage}
\tableofcontents
\clearpage
\part*{I.\quad Before Bohr: the Danish philosophical corpus}
\addcontentsline{toc}{part}{I.\quad Before Bohr}

\noindent Occurrences in the sks-search index (SKS, and the transcriptions and e-books of the other authors). \emph{entydig} and \emph{eentydig} occur 3 times in all, every one below; the old word for the opposite, \emph{tvetydig}, and its negation \emph{utvetydig} are counted for comparison.

\begin{center}\begin{tabular}{@{}lrr@{}}\toprule
author & \emph{tvetydig*} & \emph{utvetydig*}\\\midrule
kierkegaard & 355 & 4 \\
nielsen & 290 & 15 \\
brochner & 52 & 6 \\
hoeffding & 17 & 3 \\
treschow & 4 & 0 \\
moller & 2 & 0 \\
sibbern & 1 & 0 \\
kroman & 0 & 1 \\
\bottomrule\end{tabular}\end{center}

\begin{quote}
Lorentz' Fortolkning er følgende: Vi kan benytte Æteren til deri
at fastlægge et Koordinatsystem og saaledes definere absolut Hvile og
absolut Bevægelse i Forhold til dette Koordinatsystem. Vi kan maaske
ogsaa definere absolut Samtidighed for Begivenheder paa forskellige
Steder og ordne Tidspunkter i Tidsfølge paa \tlg{entydig} Maade, idet vi i
det mindste kan forestille os vilkaarlig store Hastigheder (altsaa paa
lignende Maade som i Redaktør Helge Holsts Artikel). Men de i
2. omtalte Grundhypoteser fører til, at naar man meddeler et System en
jævn Translation i Forhold til Æteren, forkortes alle Længder i
Bevægelsesretningen, og alle Ure giver sig til at gaa langsommere (de
elastiske Kræfter i f. Eks. Uroen forandrer sig jo ogsaa), saa at
ethvert System efter Lorentz faar sin effektive
Tid. Disse Ændringer, der jo er udtrykt ved de
Lorentz-Einsteinske Transformationsformler, tilfredsstiller derfor
ogsaa den af Einstein først paapegede Reciprocitet: Dersom $A$ og $B$
bevæger sig i Forhold til hinanden med en vis relativ Hastighed, saa at
de bevæger sig med forskellig Hastighed i Forhold til Æteren, eller at
en af dem maaske er i Hvile i Forhold til denne, maa $A$ mene nøjagtig
det samme om $B$ (d. v. s. $B$'s Maalestokke og Ure), som $B$ mener om
$A$. Det er derfor udelukket, at nogen af dem eller nogen anden
kan fastslaa, hvem af dem der muligvis hviler i Forhold til
Æteren, eller hvem der rigtigt kan afgøre, hvorvidt to Begivenheder paa
to forskellige Steder er samtidige eller ej [note: Lorentz' ogsaa af
Professor Kroman omtalte Diskussion mellem A og B (Das
Relativitätsprinzip, S. 22–23) skal kun oplyse dette Forhold. Hvis der
efter Lorentz bliver mindre heftig Strid mellem A og B end efter
Einstein, ligger det kun deri, at hvis de er Lorentztilhængere, kan de
enes om en fælles ideal Definition af Hvile og Samtidighed; de kan
aldrig afgøre, hvis Opfattelse der er den ideale, eller om nogen er
det. Men da A og B er Fysikere, vil de være for fornuftige til at føre
\end{quote}
\dcite{kroman/hansen-svar:14  (H. M. Hansens svar til Kroman; p. 21+) — 1916, Fysisk Tidsskrift 15 (1916–17), pp. 18–29}

\begin{quote}
S. 21 indvendes, at Tanken ikke kan fatte, at et Lyspunkt samtidig
skulde kunne være paa to Steder. Denne Indvending forstaar jeg næppe
helt. (At Transformationsligningerne, som er fundne ud fra den
Forudsætning, at Lysets Hastighed i begge Systemer i alle Retninger er
$c$, fører til det Resultat, at den er $c$, er jo ikke mærkeligt). Fra
Lorentz' Standpunkt er Sagen imidlertid simpel nok: i Følge dette kan
vi altid tænke os angivet \tlg{entydigt}, hvor i Æteren Lysblinket er naaet
hen, dets absolutte Koordinater. Blot kan hverken $A$ eller $B$ maale
andet end dets relative Koordinater, nemlig fastslaa, hvor Lysblinket
befinder sig i deres Systemer. De finder derfor hver sit Udtryk for,
hvor det er, men disse Udtryk betegner til samme absolutte Tid det
samme Sted i Æteren, det samme absolutte Sted. Einstein kender intet
Middel til at angive Lysblinkets absolutte Sted og kan som ofte nævnt
ikke forbinde nogen Mening med at tale om dets absolutte Sted. Men naar
Lysglimtet passerer $A$ (Koordinat $x$ i $A$'s System) i det Øjeblik,
Klokken dér er $t$, saa tillader Formlerne \tlg{entydigt} at udregne, hvor
$B$, som vi tænker os passere $A$ i samme Øjeblik, befinder sig i sit
System (Koordinat $x'$), og hvad han mener, Klokken dér er ($t'$). Og
han vil da mene, at Lysblinket netop til Tiden $t'$ er naaet til $x'$.
\end{quote}
\dcite{kroman/hansen-svar:30  (H. M. Hansens svar til Kroman; p. 26) — 1916, Fysisk Tidsskrift 15 (1916–17), pp. 18–29}

\begin{quote}
Efter Relativitetsteoriens Fremkomst fremhæves det imidlertid som Einsteins
Fortjeneste, at han har indset Nødvendigheden af at give en bestemt Definition
af Samtidighed mellem Begivenheder paa forskellige Steder i Rummet og tillige
gjort den mærkelige Opdagelse, at det var umuligt at gøre en saadan Definition
\tlg{eentydig}. Sandheden turde dog snarere være den, at det Samtidighedsbegreb, der
følger af vor Bevidstheds Natur, i Virkeligheden tidligere uden Definition
indgik paa rette Maade i Fysikernes Tankerækker, medens Einsteins Definition
har skabt Forvirring, fordi den har bevirket, at Begrænsninger i vore fysiske
Midler til Samtidighedsbestemmelse er blevet ophøjet til noget grundvæsentligt
for vor Erkendelse. Men da der er vakt Tvivl i denne Sag, kan det være nyttigt
at gøre sig helt klart, hvad der er det virkelige Indhold af det hidtil uden
nærmere Forklaring anvendte Samtidighedsbegreb (for Begivenheder paa
forskellige Steder), og dette er i Grunden ikke vanskeligt.
\end{quote}
\dcite{kroman/holst-tidsproblemet:5  (Tidsproblemet; p. 2+) — 1916, Fysisk Tidsskrift 15 (1916–17), pp. 1–13}

\part*{II.\quad Bohr's Danish in print}
\addcontentsline{toc}{part}{II.\quad Bohr's Danish in print}

\noindent Each passage was transcribed from the page image and checked against it; the page is the one printed on it.

\section{1925: Atomteori og Mekanik — the Danish manuscript beside the published versions}

\noindent From the Collected Works' table “Examples of variants in the versions A, B, D and F”, with the original (uncorrected) Danish manuscript of the address given at the Sixth Scandinavian Mathematics Congress, Copenhagen, 31 August 1925. A = Nature (Suppl.) 116 (1925); B = the congress Beretning (1926), English; D = Naturwissenschaften 14 (1926), German; F = Atomteori og Naturbeskrivelse (1929), Danish. The Collected Works judges A the later of the two English versions. Niels Bohr Collected Works, vol. ?, p. [271].

{\small\begin{longtable}{@{}>{\raggedright\arraybackslash}p{0.160\textwidth}>{\raggedright\arraybackslash}p{0.160\textwidth}>{\raggedright\arraybackslash}p{0.160\textwidth}>{\raggedright\arraybackslash}p{0.160\textwidth}>{\raggedright\arraybackslash}p{0.160\textwidth}@{}}\toprule
\textbf{A (English)} & \textbf{B (English)} & \textbf{D (German)} & \textbf{F (Danish, 1929)} & \textbf{Danish manuscript (1925)}\\\midrule
p. 850: which prevents a \olg{unique} assignment of quantum indices on the basis of mechanical pictures. & p. 64: [As A.] & p. 7: der eine \olg{eindeutige} Zuordnung von Quantenzahlen in Anlehnung an die Quantisierungsregeln für Periodizitätssysteme ausschließt. & p. 34: der forhindrer en \tlg{entydig} Benyttelse af Kvantetal i Tilknytning til Kvantiseringsreglerne for Periodicitetssystemer. & p. 10: som udelukker en \tlg{entydig} Tildeling af Kvantetal i Tilknytning til Periodicitetssystemernes Teori.\\[1ex]
p. 851: that the results do not allow of a \olg{unique} association with mechanical pictures. & p. 65: that the results do not allow of an association with a \olg{unique} mechanical picture. & p. 7: daß sie keine \olg{eindeutige} Zuordnung zu mechanischen Bildern zulassen. & p. 35: at de ikke tillader nogen \tlg{entydig} Benyttelse af mekaniske Billeder. & p. 10: saa unddrager disse Resultater sig for det første en \tlg{entydig} Indordning under mekaniske Forestillinger.\\[1ex]
\bottomrule\end{longtable}}
\section{1929: Atomteorien og Grundprincipperne for Naturbeskrivelsen}

\noindent Address to the 18th Meeting of Scandinavian Scientists, Copenhagen, 26 August 1929; Beretning om det 18. skandinaviske Naturforskermøde (1929), pp. 71–83. Bohr's Danish original.

\begin{quote}
Men dersom de paagældende fysiologiske Fænomener frembyder en til den omhandlede Grænse udviklet Forfining, betyder det jo, at vi samtidig nærmer os Grænsen for deres \tlg{entydige} Beskrivelse ved Hjælp af vore sædvanlige anskuelige Forestillinger.
\end{quote}
\dcite{Beretning om det 18. skandinaviske Naturforskermøde (1929), p. 82; Collected Works p. [234]}

\section{1956: Atomerne og den menneskelige erkendelse}

\noindent Separate print, København 1956 (in the folder as SP\_57\_00\_00\_1956\_1761.pdf). The same essay stands in Filosofiske Skrifter II, where it has no pages.

\begin{quote}
Omvendt kræver enhver \tlg{entydig} anvendelse af dynamiske bevarelseslove i kvantefysikken, at det omhandlede fænomens beskrivelse for de atomare objekters vedkommende indebærer et principielt afkald med hensyn til detaljeret rum-tidskoordination.
\end{quote}
\dcite{Atomerne og den menneskelige erkendelse (København 1956), p. 9}

\begin{quote}
I kvantefysikken er, som vi har set, en redegørelse for måleinstrumenternes funktioner uundværlig for definitionen af fænomenerne, og vi må så at sige drage skillelinien mellem subjekt og objekt på en måde, der i hvert enkelt tilfælde sikrer den \tlg{entydige} anvendelse af de i meddelelserne benyttede elementære fysiske begreber.
\end{quote}
\dcite{Atomerne og den menneskelige erkendelse (København 1956), p. 10}

\section{1959: letter to Max Delbrück, 25 July 1959}

\noindent Carbon copy, Tisvilde. Printed in the Collected Works, Part V, Selected Correspondence (vol. ?), pp. [479]–[480]; translation p. [481]. Note the spelling eentydig, and its pairing with tvetydighed and with "objektiv". A letter, not a publication.

\begin{quote}
Naturligvis rummer enhver brug af ordet “tænke” med dets implicite henvisning til et bevidst subjekt en tvetydighed, så længe man ikke klargør sig sprogets formål, at meddele sig på \tlg{eentydig} eller måske rettere objektiv måde til vore medmennesker.
\end{quote}
\dcite{Bohr to Delbrück, 25 July 1959; Collected Works p. [479]}

\begin{quote}
Pointen er jo her, at selvom enhver \tlg{eentydig} meddelelse kræver distinktion mellem et subjekt og et objekt, kan det subjekt, der underforstås i en given situation, helt eller delvis indgå i det objektive indhold af en meddelelse om en anden situation.
\end{quote}
\dcite{Bohr to Delbrück, 25 July 1959; Collected Works p. [480]}

\section{1961: Atomvidenskaben og menneskehedens krise}

\noindent Politiken, 20 April 1961. Printed in the Collected Works under TEXT AND TRANSLATION (vol. ?), p. [284].

\begin{quote}
Trods den gennemgribende revision af grundlaget for den \tlg{entydige} anvendelse af så elementære begreber som rum og tid, kan relativitetsteorien betragtes som en harmonisk afrunding af den såkaldte klassiske fysik.
\end{quote}
\dcite{Politiken, 20 April 1961; Collected Works p. [284]}

\begin{quote}
Derimod skulle opdagelsen af det universelle virkningskvantum afsløre et helhedstræk ved de atomare processer, som var ganske fremmed for den mekaniske naturopfattelse, der jo antager muligheden af en uindskrænket opdeling af fænomenerne og af deres beskrivelse som en kæde, hvis enkelte led ligegyldigt hvor langt underdelingen fortsættes kan sammenknyttes gennem \tlg{entydig} anvendelse af begreberne årsag og virkning.
\end{quote}
\dcite{Politiken, 20 April 1961; Collected Works p. [284]}

\part*{III.\quad Bohr's Danish collections: Filosofiske Skrifter I--II}
\addcontentsline{toc}{part}{III.\quad Filosofiske Skrifter I--II}

\noindent Every paragraph of \emph{Filosofiske Skrifter} I and II with \emph{entydig}. The Danish is unpaginated there. Beneath each, in small type, is the paragraph of the printed original it was aligned to (OCR of the \emph{Collected Works} extract), with the words that correspond marked \olg{like this}. An alignment marked ``$\approx$'' is right about nine times in ten; the usual error is a slip of one paragraph.

\section{Filosofiske Skrifter I}

\subsection{Indledende oversigt}

\begin{quote}
I den tredie artikel, der er et bidrag til et af tidsskriftet »Die Naturwissenschaften« i anledning af Plancks 50-års doktorjubilæum i juni 1929 udgivet festskrift, er der mere udførligt gået ind på kvanteteoriens almindelige filosofiske side. Ikke mindst med henblik på den beklagelse, som i vide kredse er kommet til orde over for afkaldet på en streng årsagsbeskrivelse af atomfænomenerne, søger forfatteren at vise, at de vanskeligheder for vore anskuelsesformer, som vi på grund af virkningskvantets udelelighed møder i atomteorien, tør betragtes som en lærerig påmindelse om de almindelige vilkår, der gælder for menneskelige begrebsdannelser. Umuligheden af på tilvant måde at skelne mellem de fysiske fænomener og deres iagttagelse stiller os i virkeligheden over for en ganske lignende situation som den, vi kender fra psykologien, hvor vi stadig påmindes om vanskeligheden ved at skille mellem subjekt og objekt. Det kunne måske ved første øjekast se ud, som om en sådan indstilling over for fysikken giver plads for en mystik, der er i modsætning til naturvidenskabens ånd. Lige så lidt som ved andre menneskelige spørgsmålstillinger turde dog klarhed på det omhandlede område kunne nås uden ved at se de vanskeligheder i øjnene, som frembyder sig for begrebsbygningen og udtryksmidlernes anvendelse. Efter forfatterens opfattelse ville det således være en misforståelse at mene, at vanskelighederne på atomteoriens område kunne omgås ved at erstatte den klassiske fysiks begreber med eventuelle nye begrebsdannelser. Som allerede fremhævet, betyder jo erkendelsen af vore anskuelsesformers begrænsning på ingen måde, at vi ved indordningen af sanseindtrykkene kan undvære de tilvante forestillinger eller deres umiddelbare udtryk i sproget. Lige så lidt turde de grundbegreber, som de klassiske fysiske teorier har skænket os, nogensinde blive overflødige for beskrivelsen af de fysiske erfaringer. Ikke alene beror erkendelsen af virkningskvantets udelelighed og bestemmelsen af dets værdi på en analyse af målinger baseret på klassiske begreber, men det er stadig udelukkende anvendelsen af disse begreber, der betinger forbindelsen mellem kvanteteoriens symbolik og erfaringernes indhold. Samtidig må vi imidlertid betænke, at muligheden for disse grundbegrebers \tlg{entydige} brug alene hviler på den indre sammenhæng hos de klassiske teorier, hvorfra de er hentede, og at grænserne for disse begrebers anvendelse derfor ifølge sagens art er betinget af det omfang, i hvilket vi ved redegørelsen for fænomenerne kan se bort fra det for de klassiske teorier fremmede element, der symboliseres ved virkningskvantet.
\end{quote}
\dcite{Filosofiske Skrifter 1, Indledende oversigt, ¶14 — no printed source indexed}

\begin{quote}
Det er netop denne sagernes stilling, der holdes os for øje ved det ofte berørte dilemma vedrørende lysets og materiens egenskaber. Kun i direkte tilknytning til den klassiske elektromagnetiske teori kan der blive tale om at give spørgsmålet om lysets og materiens væsen et håndgribeligt indhold. Vel er lyskvanter og materiebølger uvurderlige hjælpemidler ved formuleringen af de statistiske lovmæssigheder, der behersker sådanne fænomener som de fotoelektriske virkninger og elektronstrålers interferens. Men ved disse fænomer befinder vi os jo netop på et område, hvor en hensyntagen til virkningskvantet er uundgåelig, og hvor en \tlg{entydig} beskrivelse er uigennemførlig. Den i denne forstand symbolske karakter af de nævnte hjælpemidler fremtræder også deri, at en udtømmende beskrivelse af de elektromagnetiske bølgefelter ikke levner nogen plads for lyskvanter, samt at der ved benyttelsen af forestillingen om materiebølgerne aldrig er tale om en fuldstændig beskrivelse i lighed med de klassiske teoriers. Som det er fremhævet i den anden artikel, kommer den absolute værdi af bølgernes såkaldte fase jo aldrig i betragtning ved tolkningen af erfaringerne. I denne forbindelse må det også betones, at betegnelsen »sandsynlighedsamplituder« for materiebølgernes amplitudefunktioner er led i en ofte bekvem udtryksmåde, der dog ikke kan gøre krav på nogen almindelig gyldighed. Som nævnt er det kun ved hjælp af de klassiske forestillinger muligt at tilskrive iagttagelsesresultaterne et \tlg{entydigt} indhold, og ved sandsynlighedsbetragtninger vil det derfor altid dreje sig om udfaldet af forsøg, der tydes ved hjælp af sådanne forestillinger. Som følge heraf vil den brug, der gøres af de symbolske hjælpemidler i hvert enkelt tilfælde afhænge af de nærmere omstændigheder vedrørende forsøgenes indretning. Det, der giver den kvanteteoretiske beskrivelse dens særpræg, er nu netop, at vi for at komme udenom virkningskvantet må benytte forskellige forsøgsindretninger for at opnå nøjagtige målinger af de forskellige størrelser, hvis samtidige kendskab ville kræves for en fuldstændig beskrivelse baseret på de klassiske teorier, samt at disse målingsresultater ikke kan suppleres ved gentagne målinger. Virkningskvantets udelelighed forlanger nemlig, at der ved tydningen af hvert enkelt målingsresultat i tilknytning til de klassiske forestillinger tillades et spillerum i vort regnskab med vekselvirkningen mellem genstand og målemiddel, der medfører, at en efterfølgende måling i et vist omfang berøver de oplysninger, som en forudgående måling har givet os, deres betydning for forudsigelser af fænomenernes fremtidige forløb. Dette forhold sætter åbenbart ikke alene en grænse for omfanget af de oplysninger, som målinger kan give os, men også for den mening vi kan tillægge sådanne oplysninger. Vi møder her i ny belysning den gamle erkendelse, at der ved naturbeskrivelsen ikke er tale om at afdække fænomernes egentlige væsen, men kun om i størst muligt omfang at efterspore sammenhæng i vore erfaringers mangfoldighed.
\end{quote}
\dcite{Filosofiske Skrifter 1, Indledende oversigt, ¶15 — no printed source indexed}

\begin{quote}
Som følge af denne situation taber jo selv ord som være og vide deres \tlg{entydige} mening. Et interessant eksempel på vor sprogbrugs tvetydighed i den omhandlede forbindelse er den talemåde, hvorefter årsagsbeskrivelsens svigten udtrykkes ved at sige, at det drejer sig om frit valg fra naturens side. Egentlig kræver jo en sådan talemåde forestillingen om en udenforstående vælger, hvad der dog allerede modsiges gennem brugen af ordet natur. Vi stilles her over for et grundtræk i det almindelige erkendelsesproblem, og vi må gøre os klart, at vi efter sagens væsen til syvende og sidst altid er henvist til at udtrykke os gennem et maleri med ord, der benyttes på uanalyseret måde. Som det betones i artiklen, må vi jo på alle erkendelsesområder erindre, at vor bevidstheds væsen betinger et komplementaritetsforhold mellem analysen af ethvert begreb og dets umiddelbare anvendelse.
\end{quote}
\dcite{Filosofiske Skrifter 1, Indledende oversigt, ¶17 — no printed source indexed}

\subsection{Atomteori og mekanik \normalfont(Nature (Suppl.) 116 (1925), pp. 845–852)}

\begin{quote}
Analysen af spektrenes finere struktur skulle imidlertid bringe et antal enkeltheder for dagen, der ikke lader sig tyde ved hjælp af mekaniske billeder på grundlag af kvantiseringsreglerne for periodicitetssystemer. Her tænkes særlig på spektralliniernes multipletstruktur og virkningerne af magnetiske kræfter på disse strukturer. De sidste virkninger, som sædvanligt betegnes som anomale Zeeman-effekter og som allerede beredte den klassiske teori vanskeligheder, indføjer sig dog utvungent i et til kvanteteoriens grundpostulater svarende skema, idet, som Landé viste, liniernes opspaltningskomponenter lader sig fremstille som kombinationer mellem de i magnetfeltet opspaltede spektralled. Sterns og Gerlachs smukke forsøg, hvorved en umiddelbar forbindelse påvistes mellem den kraft, der angriber et atom i et inhomogent magnetisk felt, og de energiværdier som ud fra de opspaltede spektralled beregnes for atomets stationære tilstande i feltet, tør endda betragtes som en hovedstøtte for kvanteteoriens grundforestillinger. Imidlertid skulle Landés analyse samtidig åbenbare en principiel forskel mellem elektronernes vekselspil i atomet og koblingen i mekaniske systemer. Vi kan sige, at elektronernes vekselspil her åbenbarer en ikke-mekanisk beskrivelig »tvang«, der forhindrer en \tlg{entydig} benyttelse af kvantetal i tilknytning til kvantiseringsreglerne for periodicitetssystemer. Af særlig betydning i forbindelse med diskussionen af disse forhold var en af Ehrenfest udledt almindelig betingelse for den termodynamiske stabilitet, der, når den anvendes sammen med kvanteteoriens postulater, udsiger, at en stationær tilstands statistiske vægt ikke ændres ved en kontinuert transformation af atomsystemets mekaniske beskaffenhed. Iøvrigt fører denne betingelse allerede ved atomer med kun en elektron til vanskeligheder, der antyder en begrænset gyldighed af teorien for periodicitetssystemer. Det mekaniske problem om atomdelenes bevægelser indeholder nemlig singulære løsninger, der må udelukkes fra de stationære tilstandes mangfoldighed. Denne udelukkelse betyder en kunstig begrænsning af kvantiseringsreglerne, der dog til at begynde med ikke stod i åbenbar modsætning til det eksperimentelle materiale. Vanskeligheder af særlig alvorlig natur er imidlertid for nylig bragt for dagen ved den af Klein og Lenz udførte interessante analyse af problemet om brintatomet i krydsede elektriske og magnetiske felter. Her er det umuligt at tilfredsstille Ehrenfests betingelse, da bevægelser, der ikke lader sig udelukke fra de stationære tilstandes mangfoldighed, ved en passende ændring af de ydre kræfter overføres i baneformer, hvor elektronen falder ind i kernen.
\end{quote}
\dcite{Filosofiske Skrifter 1, Atomteori og mekanik, ¶23 — ≈ pp. 850–851 — aligned to the printed text, not read off a page; that page is itself inferred}
{\footnotesize\begin{quote}\itshape
I n fact, we are forced t o assume the presence of a mechanically undescribable ‘ I strain ” in the interaction of the electrons which prevents a \olg{unique} assignment of quantum indices on the basis of mechanical p i c t \textasciitilde{} r e s .
\end{quote}}

\begin{quote}
Til trods for de nævnte vanskeligheder skulle analysen af spektrenes finere enkeltheder alligevel væsentlig fremme den kvanteteoretiske fortolkning af lovene for grundstoffernes slægtskabsforhold. Med henvisning til en række kendsgerninger af forskellig art har således Dauvillier, Main Smith og Stoner gjort betydningsfulde forslag til den nærmere udformning af kvanteteoriens almindelige forestillinger vedrørende elektronernes gruppeinddeling i atomet. Disse betragtninger udviser uanset deres formelle natur en nøje forbindelse med de spektrale lovmæssigheder, der blev opdaget ved Landés analyse. Ad denne vej er der fornylig især af Pauli opnået meget lovende fremskridt. Omend disse resultater betyder et vigtigt skridt til opfyldelsen af det nævnte program om forklaringen af grundstoffernes egenskaber alene ved hjælp af atomnummeret, må man ikke glemme, at de ikke tillader nogen \tlg{entydig} benyttelse af mekaniske billeder.
\end{quote}
\dcite{Filosofiske Skrifter 1, Atomteori og mekanik, ¶24 — ≈ p. 851 — aligned to the printed text, not read off a page; that page is itself inferred}
{\footnotesize\begin{quote}\itshape
In this direction important progress has recently been made, especially by Pauli.9 Sotwithstanding that the results thus obtained constitute an important step towards the above-mentioned programme of accounting for the properties of the elements solely on the basis of the atomic number, it must be remembered, however, that the results do not allow of a \olg{unique} association with mechanical pictures.
\end{quote}}

\subsection{Kvantepostulatet og atomteoriens seneste udvikling \normalfont(Atti del Congresso Internazionale dei Fisici, 11–20 Settembre 1927, Como–Pavia–Roma, Volume Secondo (Nicola Zanichelli, Bologna 1928), pp. 565–588 — first version)}

\begin{quote}
Dette forhold indebærer vidtgående konsekvenser. På den ene side fordrer definitionen af et fysisk systems tilstand, som man sædvanlig opfatter det, udelukkelsen af alle ydre indvirkninger; men så er også ifølge kvantepostulatet enhver mulighed for iagttagelse udelukket, og fremfor alt taber begreberne tid og rum deres umiddelbare betydning. Tillader vi på den anden side for at muliggøre iagttagelse eventuelle vekselvirkninger med dertil egnede ikke til systemet hørende målemidler, er ifølge sagens natur en \tlg{entydig} definition af systemets tilstand ikke mere mulig, og der kan ikke blive tale om kausalitet i sædvanlig forstand. Ifølge kvanteteoriens væsen må vi altså nøjes med at opfatte rum-tidsbeskrivelsen og kausalitetsfordringen, hvis forening karakteriserer de klassiske teorier og som symboliserer iagttagelses-og definitionsmulighedernes idealisation, som komplementære, men hinanden udelukkende træk i beskrivelsen af erfaringens indhold. Ligesom man i relativitetsteorien erkender, at hensigtsmæssigheden af den skarpe, af vore sanser krævede adskillelse mellem rum og tid kun beror derpå, at de sædvanligt forekommende relative hastigheder er små i forhold til lysets hastighed, kan kvanteteorien siges at have bragt den erkendelse, at formålstjenligheden af vor hele kausale rum-tidslige anskuelsesmåde kun beror på virkningskvantets lidenhed i forhold til de virkninger, der kommer i betragtning ved de sædvanlige sansefornemmelser. I virkeligheden stiller kvantepostulatet os ved beskrivelsen af de atomare fænomener over for den opgave at udvikle en »komplementaritetsteori«, hvis modsigelsesfrihed kun kan bedømmes ved at veje definitions- og iagttagelsesmulighederne mod hverandre.
\end{quote}
\dcite{Filosofiske Skrifter 1, Kvantepostulatet og atomteoriens seneste udvikling, ¶4 — ≈ pp. 566–567 — aligned to the printed text, not read off a page; that page is itself inferred}
{\footnotesize\begin{quote}\itshape
On the other hand, i f in order to make observation possible we permit certain interactions with suitable means of measurement, not belonging to the system, an \olg{unambiguous} definition of its state is naturally no longer possible, and there can be no question of causality in the ordinary sense of the word.
\end{quote}}

\begin{quote}
for den størst mulige definitionsnøjagtighed af de bølgefeltet tilordnede individers energi og impuls. I almindelighed vil dog forholdene for tilordningen af en energi- og impulsværdi til et bølgefelt være endnu mindre gunstige. Selv om bølgegruppens beskaffenhed til at begynde med opfylder relationen (2), vil dens udstrækning i tidens løb undergå sådanne forandringer, at den bliver mindre og mindre egnet til at fremstille et individuums bevægelse. Netop i denne omstændighed ligger jo det paradoksale i spørgsmålet om lysets og de materielle smådeles natur. Iøvrigt hænger den ved formel (2) udtrykte begrænsning af de klassiske begreber nøje sammen med den indskrænkede gyldighed af den klassiske mekanik, som i materiens bølgeteori svarer til den geometriske optik, i hvilken bølgeudbredelsen anskueliggøres ved »stråler«. Kun i dette grænsetilfælde lader energi og impuls sig \tlg{entydigt} definere i tilslutning til rum-tidsbilleder. For en almindelig definition af disse begreber er vi direkte henvist til sætningerne om energiens og bevægelsesmængdens bevarelse, hvis rette formulering har været et grundproblem ved udviklingen af de symbolske metoder, vi senere skal komme ind på.
\end{quote}
\dcite{Filosofiske Skrifter 1, Kvantepostulatet og atomteoriens seneste udvikling, ¶14 — the paper is known, the paragraph is not — no page}

\begin{quote}
Ligesom ved stedbestemmelsen kan iagttagelsesprocessens varighed ved impulsmålinger gøres vilkårlig kort, hvis vi blot benytter tilstrækkelig kortbølget stråling. At tilbagestødet derved bliver større, har jo, som vi har set, ingen indflydelse på målenøjagtigheden. Det må endnu bemærkes, at det, når vi gentagne gange har talt om partiklens hastighed, i denne sammenhæng kun drejer sig om en bekvem, formel tilslutning til den sædvanlige rum-tidsbeskrivelse. Som det allerede fremgår af de ovenanførte betragtninger af de Broglie, må hastighedsbegrebet stedse anvendes med forsigtighed. En \tlg{entydig} definition af dette begreb er jo også udelukket ved kvantepostulatet, hvad man særlig må betænke, når man sammenligner resultatet af flere på hinanden følgende iagttagelser. Vel kan man bestemme et individs sted til to forskellige tidspunkter med en hvilken som helst nøjagtighed. Hvis vi imidlertid på sædvanlig måde vil beregne individets hastighed i intervallet derimellem, så har vi at gøre med en idealisation, ud fra hvilken vi ikke kan drage \tlg{entydige} slutninger med hensyn til individets tidligere og fremtidige forhold.
\end{quote}
\dcite{Filosofiske Skrifter 1, Kvantepostulatet og atomteoriens seneste udvikling, ¶22 — ≈ pp. 574–575 — aligned to the printed text, not read off a page}
{\footnotesize\begin{quote}\itshape
The purpose, however. \textasciitilde{} 7 a s only to obtain a connection w-ith our ordinary space-time description con>-enient in this case, since, strictly speaking, an \olg{unambiguous} definition of velocity is excluded by the quantum postulate. … We haire seen that the position of an individual at two given moments can be determined to any desired degree of accuracy; but if we from such measurements calculate the 1-elocity of the individual in the ordinary way, it must he clearly realized that we ase dealing with an abstraction, from which no \olg{unambiguous} information concerning the previous or future behaviour of the individual can be obtained. -4s shown above, this results from the fact that the determination of the position of a particle always involves a complete rupture in the description of its dynamical behaviour, in the same way as the accurate determination of its momentum always involves a gap in the knowledge of its spatial propagation.
\end{quote}}

\begin{quote}
Når den Schrödingerske bølgelignings egensvingninger giver en velegnet fremstilling af atomets stationære tilstande, er det netop fordi de i forbindelse med den almindelige kvanterelation (I) tillader en \tlg{entydig} definition af systemets energi. Dette kræver imidlertid en vidtgående resignation med hensyn til rum-tidsbeskrivelsen, der, som vi skal se, udelukker enhver mulighed for at forbinde et nærmere kendskab til de enkelte partiklers forhold i atomet med en modsigelsesfri anvendelse af begrebet stationære tilstande. Ved problemer, hvor en forfølgelse af atomdelene i tid og rum indgår som et væsentligt led ved iagttagelsernes tydning, er vi henvist til en undersøgelse af bølgeligningens almindelige løsning, som fremkommer ved en superposition af egenløsninger. Det drejer sig her om en komplementaritet af definitionsmulighederne af ganske samme art som ved det tidligere betragtede spørgsmål om lysets og de frie materielle partiklers egenskaber. Medens definitionen af individernes energi og impuls var knyttet til begrebet harmonisk elementarbølge, hvilede, som vi så, enhver beskrivelse af fænomenerne i rum og tid på en betragtning af de interferenser, som finder sted inden for en gruppe af sådanne elementarbølger. Også i det her omhandlede tilfælde lader overensstemmelsen mellem iagttagelsesmulighederne og definitionsmulighederne sig direkte eftervise.
\end{quote}
\dcite{Filosofiske Skrifter 1, Kvantepostulatet og atomteoriens seneste udvikling, ¶40 — ≈ pp. 586–587 — aligned to the printed text, not read off a page; that page is itself inferred}
{\footnotesize\begin{quote}\itshape (No corresponding word in the aligned paragraph of the original.)\end{quote}}

\begin{quote}
Ifølge kvantepostulatet vil en iagttagelse, der kan belære os om de enkelte atomdeles forhold, altid medføre en ændring af atomets tilstand. Som fremhævet af Heisenberg, vil denne ændring, når det gælder atomer i stationære tilstande med lave kvantetal, endogså i almindelighed bestå i en udstødning af den betragtede elektron af atomet. En bestemmelse af elektronens »bane« i atomet ved gentagne iagttagelser er altså i sådanne tilfælde udelukket. Dette hænger sammen med den omstændighed, at man af egensvingninger med få knuder ikke kan opbygge en bølgegruppe, som blot tilnærmelsesvis kan fremstille en partikels »bevægelse«. Beskrivelsens komplementære natur kommer imidlertid fremfor alt til udtryk derved, at en \tlg{entydig} fortolkning af iagttagelser over partiklernes forhold i atomet altid må bero på, at man under iagttagelsesprocessen kan se bort fra vekselvirkningen og betragte partiklerne som frie. Dette fordrer imidlertid, at processens varighed er kort i forhold til atomets naturlige perioder og medfører derfor en usikkerhed i vort kendskab til den iagttagelsen ledsagende energiændring, som er stor i forhold til energidifferenserne mellem på hinanden følgende stationære tilstande.
\end{quote}
\dcite{Filosofiske Skrifter 1, Kvantepostulatet og atomteoriens seneste udvikling, ¶41 — ≈ p. 587 — aligned to the printed text, not read off a page; that page is itself inferred}
{\footnotesize\begin{quote}\itshape (No corresponding word in the aligned paragraph of the original.)\end{quote}}

\begin{quote}
Ved begrebet stationære tilstande har vi, som allerede nævnt, at gøre med en karakteristisk anvendelse af kvantepostulatet. Dette begreb fordrer ifølge sin natur, at der ses fuldstændig bort fra en tidsbeskrivelse, hvad der jo ud fra det her indtagne synspunkt netop også er betingelsen for en \tlg{entydig} definition af atomets energi. Strengt taget fordrer begrebet stationær tilstand udelukkelsen af enhver ydre vekselvirkning med individer der ikke hører til systemet. At der tilskrives et sådant afsluttet system en bestemt energiværdi, kan opfattes som et umiddelbart udtryk for den i energisætningen indeholdte kausalitetsfordring. Heri ser vi berettigelsen af den til grund for anvendelsen af kvantepostulatet på spørgsmålet om atomernes bygning liggende antagelse om de stationære tilstandes supramekaniske stabilitet, ifølge hvilken atomet før og efter enhver ydre indvirkning befinder sig i en veldefineret stationær tilstand.
\end{quote}
\dcite{Filosofiske Skrifter 1, Kvantepostulatet og atomteoriens seneste udvikling, ¶44 — ≈ p. 589 — aligned to the printed text, not read off a page; that page is itself inferred}
{\footnotesize\begin{quote}\itshape (No corresponding word in the aligned paragraph of the original.)\end{quote}}

\begin{quote}
De hidtidige metoders utilstrækkelighed overfor elementarpartiklernes problem kommer ved de lige omtalte spørgsmål frem derved, at de ikke tillader en \tlg{entydig} begrundelse af det af Pauli opstillede såkaldte udelukkelsesprincip, der åbenbarer så ejendommelig en forskel mellem de elektriske elementarpartikler og de af lyskvanteforestillingen symboliserede »individer«. I dette princip, som har været overordentlig frugtbart såvel for problemet om atomernes bygning som for den nyeste udvikling af de statistiske teorier, har vi jo at gøre med en af flere tænkelige muligheder, som hver for sig tilfredsstiller korrespondensfordringerne, iøvrigt møder vi ved magnetelektronproblemet et særlig lærerigt eksempel på den vanskelighed, det volder i kvanteteorien at tilfredsstille relativitetsfordringen. Således var det hidtil ikke muligt at bringe Darwins og Paulis iøvrigt så lovende forsøg på en til behandling af dette problem egnet almindeliggørelse af de kvanteteoretiske metoder i overensstemmelse med Thomas’ forklaring af de eksperimentelle resultater så væsentlige relativitetskinematiske betragtninger. I den allerseneste tid er det imidlertid lykkedes Dirac med held at angribe problemet om den magnetiske elektron ved hjælp af en ny, yderst sindrig udvidelse af den symbolske metode, der i alle enkeltheder gør rede for de pågældende spektrale fænomener og samtidig tilfredsstiller relativitetsfordringen. Dette angreb hviler ikke blot som de tidligere fremgangsmåder på indførelsen af en kompleksitet, der er kendetegnet ved brugen af imaginære størrelser, men benytter i selve grundligningerne tallegemer af en endnu højere kompleksitetsgrad.
\end{quote}
\dcite{Filosofiske Skrifter 1, Kvantepostulatet og atomteoriens seneste udvikling, ¶55 — ≈ p. 597 — aligned to the printed text, not read off a page}
{\footnotesize\begin{quote}\itshape (No corresponding word in the aligned paragraph of the original.)\end{quote}}

\subsection{Virkningskvantet og naturbeskrivelsen \normalfont(Naturwissenschaften 17 (1929), pp. 483–486)}

\begin{quote}
Under disse forhold kan vi ikke undre os over, at det ved alle følgerigtige anvendelser af kvanteteorien stedse har drejet sig om væsentligt statistiske problemer. I Plancks oprindelige arbejder var det jo i første linie nødvendigheden af en ændring af den klassiske statistiske mekanik, som var anledningen til virkningskvantets indførelse. Dette for kvanteteorien ejendommelige træk kommer på slående måde til udtryk i den fornyede diskussion om lysets og de materielle elementarpartiklers natur. Medens disse spørgsmål tilsyneladende havde fundet deres endelige løsning inden for de klassiske teoriers ramme, ved vi nu, at vi såvel for lysets som for de materielle smådeles vedkommende behøver forskelligartede billeder for at opnå et alsidigt udtryk for fænomenerne og en \tlg{entydig} formulering af de statistiske love, som behersker iagttagelsesresultaterne. Jo klarere umuligheden af en ensartet formulering af kvanteteoriens indhold ved hjælp af klassiske forestillinger fremtræder, desto mere beundrer vi Plancks lykkelige intuition ved valget af betegnelsen virkningskvantum, der viser direkte hen til en svigten af virkningsprin-cippet, hvis centrale stilling i den klassiske naturbeskrivelse han selv ved mange lejligheder har betonet. Dette princip symboliserer så at sige det ejendommelige reciprokke symmetriske forhold mellem rum-tidsbeskrivelsen og sætningerne om energiens og bevægelsesmængdens bevarelse, hvis store frugtbarhed allerede i den klassiske fysik hænger sammen dermed, at disse sætninger i vid udstrækning kan anvendes uafhængigt af fænomenernes forfølgelse i rum og tid. Det er netop denne reciprocitet, der på heldigste måde er udnyttet i kvantemekanikkens formalisme. Faktisk optræder virkningskvantet her kun i relationer, i hvilke de i Hamiltonsk forstand kanonisk konjugerede rum-tidskoordinater og impuls-energikomponenter indgår på symmetrisk og reciprok måde. Også analogien mellem optik og mekanik, som har vist sig så frugtbar for kvanteteoriens nyeste udvikling, hænger på nøjeste måde sammen med disse forhold.
\end{quote}
\dcite{Filosofiske Skrifter 1, Virkningskvantet og naturbeskrivelsen, ¶3 — ≈ p. 484 — aligned to the printed text, not read off a page}
{\footnotesize\begin{quote}\itshape
Während diese Fragen im Rahmen der klassischen Theorien eine scheinbarendgültigeLösung bekommen hatten, so wissen wir jetzt, daß sowohl für das Licht wie für die materiellen Teilchen verschiedenartige Bilder notwendig sind, um die Erscheinungen allseitig zum Ausdruck zu bringen und eine \olg{eindeutige} Formulierung der statistischen Gesetze, welche die Beobachtungsergebnisse regeln, zu gewähren.
\end{quote}}

\begin{quote}
Det erkendelsesproblem, som det drejer sig om, er i korthed følgende: en beskrivelse af vor tankevirksomhed forlanger på den ene side, at et objektivt givet indhold stilles over for et betragtende subjekt, medens på den anden side – således som det fremgår allerede af et sådant udsagn – ingen skarp adskillelse mellem objekt og subjekt kan opretholdes, da jo også det sidstnævnte begreb tilhører vort tankeindhold. Af disse omstændigheder følger ikke alene den relative, af det vilkårlige valg af vort synspunkt afhængige betydning af ethvert begreb eller snarere ethvert ord; men vi må i almindelighed være forberedt på, at en alsidig belysning af en og samme genstand kan forlange forskellige synspunkter, der forhindrer en \tlg{entydig} beskrivelse. Strengt taget står jo den bevidste analyse af ethvert begreb i et udelukkelsesforhold til dets umiddelbare anvendelse. Med nødvendigheden af at tage vor tilflugt til en i denne forstand komplementær eller snarere reciprok beskrivelsesmåde er vi vel navnlig fortrolige gennem psykologiske problemer. Derimod turde sædvanligvis det ejendommelige for de såkaldte eksakte videnskaber ses i bestræbelsen på at opnå \tlg{entydighed} ved at undgå enhver henvisning til det betragtende subjekt. Denne bestræbelse møder vi måske mest bevidst i den matematiske symbolik, der holder os et ideal af objektivitet for øje, for hvis opnåelse der næppe er sat grænser, sålænge vi holder os inden for et i sig selv afsluttet anvendelsesområde af logikken. I de egentlige naturvidenskaber kan der imidlertid ikke være tale om noget strengt afsluttet anvendelsesområde for de logiske principper, da vi altid må regne med nye tilkommende kendsgerninger, hvis indordning i de tidligere erfaringers ramme kan kræve en revision af vore grundbegreber.
\end{quote}
\dcite{Filosofiske Skrifter 1, Virkningskvantet og naturbeskrivelsen, ¶5 — ≈ p. 485 — aligned to the printed text, not read off a page}
{\footnotesize\begin{quote}\itshape (No corresponding word in the aligned paragraph of the original.)\end{quote}}

\begin{quote}
En sådan revision har vi nylig oplevet ved skabelsen af relativitetsteorien, som netop gennem en vidtgående uddybning af iagttagelsesproblemet skulle åbenbare den subjektive karakter af alle den klassiske fysiks begreber. Trods de store fordringer, som den stiller til vor abstraktionsevne, kommer relativitetsteorien alligevel det klassiske ideal af enhed og årsagssammenhæng i naturbeskrivelsen i særlig høj grad i møde. Fremfor alt bliver forestillingen om den objektive realitet af de til vor iagttagelse kommende fænomener endnu strengt opretholdt. Som Einstein har betonet, er det jo en for hele relativitetsteorien grundlæggende antagelse, at hver iagttagelse til syvende og sidst beror på en sammentræffen af genstand og målemiddel i det samme rumtidspunkt, og forsåvidt er uafhængig af iagttagerens henførelsessystem. Efter virkningskvantets opdagelse ved vi imidlertid, at det klassiske ideal ikke kan nås ved beskrivelsen af de atomare fænomener. I særdeleshed medfører ethvert forsøg på en rum-tidsindordning et brud i årsagskæden, idet det er forbundet med en omsætning af bevægelsesmængde og energi mellem individerne og de til sammenligning benyttede målestokke og ure, fra hvilken vi ikke kan se bort, og som vi dog ikke kan føre regnskab med, når målemidlerne skal opfylde deres hensigt. Omvendt forlanger enhver \tlg{entydigt} på streng bevarelse af energi og bevægelsesmængde grundet slutning om individernes dynamiske forhold åbenbart en fuldstændig resignation med hensyn til deres forfølgelse i rum og tid. Overhovedet kan vi sige, at den kausale rum-tidsbeskrivelses hensigtsmæssighed ved de sædvanlige erfaringers indordning udelukkende beror på virkningskvantets lidenhed i sammenligning med de virkninger, der kommer i betragtning ved de sædvanlige fænomener. Plancks opdagelse har her stillet os over for en lignende situation som den, opdagelsen af lysets endelige hastighed havde bragt; hensigtsmæssigheden af den skarpe, af vore sanser forlangte adskillelse mellem tid og rum beror jo udelukkende på lidenheden af de hastigheder, som vi i det daglige liv har at gøre med, sammenlignet med lyshastigheden. I virkeligheden må ved spørgsmålet om de atomare fænomeners kausalitet målingsresultaternes reciprocitet lige så lidt glemmes som deres relativitet ved spørgsmålet om samtidigheden.
\end{quote}
\dcite{Filosofiske Skrifter 1, Virkningskvantet og naturbeskrivelsen, ¶6 — ≈ p. 486 — aligned to the printed text, not read off a page}
{\footnotesize\begin{quote}\itshape
So ist öfters die Ansicht vertreten worden, daß eine wohl nicht ausführbare, aber doch denkbare, ins einzelne gehende Verfolgung der Gehirnprozesse eine Ursachskette entschleiern würde, die eine \olg{eindeutige} Abbildung des gefühlsbetonten psychischen Geschehens darbieten würde.
\end{quote}}

\begin{quote}
Det må måske her være tilladt endnu kort at henvise til det forhold, som består mellem lovmæssighederne på det psykiske område og spørgsmålet om de fysiske fænomeners kausalitet. Med henblik på modsætningen mellem den følelse af viljesfrihed, som behersker sjælelivet, og den tilsyneladende strenge årsagssammenhæng, som de ledsagende, fysiologiske processer udviser, er den tanke jo ikke undgået filosofferne, at det her kan dreje sig om et uanskueligt komplementaritetsforhold. Den opfattelse er ofte blevet hævdet, at en vel ikke gennemførlig, men dog tænkelig, i det enkelte gående undersøgelse af hjerneprocesserne ville åbenbare en årsagskæde, der frembød en \tlg{entydig} afbildning af de følelsesbetonede sjælsoplevelser. Et sådant tankeeksperiment kommer imidlertid i ny belysning, da vi efter virkningskvantets opdagelse har lært, at en i enkeltheder gående kausal forfølgelse af atomprocesser ikke er mulig, og at ethvert forsøg på at opnå kendskab til sådanne processer betyder en principielt ukontrollerbar indgriben i deres forløb. Efter den omtalte opfattelse af forholdet mellem hjerneprocesserne og sjæls-oplevelserne må vi altså være forberedt på, at et forsøg på at iagttage de førstnævnte vil medføre en væsentlig ændring af viljesfølelsen. Om det end foreløbig her kun kan dreje sig om mere eller mindre træffende analogier, turde man dog vanskeligt kunne frigøre sig fra den overbevisning, at vi i de af kvanteteorien afslørede kendsgerninger, der ligger uden for vore sædvanlige anskuelsesformers område, har fået et middel i hænde til belysning af almindelige, filosofiske problemer.
\end{quote}
\dcite{Filosofiske Skrifter 1, Virkningskvantet og naturbeskrivelsen, ¶8 — ≈ pp. 96–97 — aligned to the printed text, not read off a page}
{\footnotesize\begin{quote}\itshape
From these circumstances follows not only the relative meaning of every concept, or rather of every word, the meaning depending upon our arbitrary choice of view point, but also that we must, in general, be prepared to accept the fact that a complete elucidation of one and the same object may require diverse points of view which defy a \olg{unique} description. … I n opposition to this, the feature which characterizes the so-called exact sciences is, in general, the attempt to attain to \olg{uniqueness} by avoiding all reference to the perceiving subject.
\end{quote}}

\subsection{Atomteorien og grundprincipperne for naturbeskrivelsen \normalfont(Beretning om det 18. skandinaviske Naturforskermøde i København, 26.–31. August 1929 (Frederiksberg Bogtrykkeri, Copenhagen 1929), pp. 71–83)}

\begin{quote}
Før jeg slutter, ligger det nær ved et sådant fællesmøde af naturforskere at berøre spørgsmålet om, hvad den beskrevne seneste udvikling af vort kendskab til atomfænomenerne kan lære os om de problemer, som de levende organismer frembyder. Selv om det vel endnu ikke er muligt at give noget fyldigt svar på dette spørgsmål, turde vi allerede skimte en vis sammenhæng mellem disse problemer og kvanteteoriens forestillingskreds. Et første fingerpeg i den retning finder vi i den omstændighed, at den til grund for sanseindtrykkene liggende vekselvirkning mellem organismerne og omverdenen i det mindste under omstændigheder kan blive så ringe, at vi nærmer os virkningskvantet. Som det ofte er bemærket, er således ganske få lyskvanter tilstrækkelige til at frembringe synsindtryk. Vi ser altså, at organismens behov, hvad selvstændighed og følsomhed angår, her er tilfredsstillet til den yderste med naturlovene forenelige grænse, og vi må være forberedt på at møde samme forhold også på andre punkter af afgørende betydning for den biologiske problemstilling. Men dersom de pågældende fysiologiske fænomener frembyder en til den omhandlede grænse udviklet forfining, betyder det jo, at vi samtidig nærmer os grænsen for deres \tlg{entydige} beskrivelse ved hjælp af vore sædvanlige anskuelige forestillinger. Dette står ingenlunde i modstrid til den kendsgerning, at de levende organismer i udstrakt grad stiller os problemer, der falder inden for vore anskuelsesformers rækkevidde, og som har afgivet så frugtbart et anvendelsesområde for fysiske og kemiske synspunkter. Heller ikke ser vi nogen umiddelbar grænse for disse synspunkters anvendelse. Lige så lidt som vi i princippet kan skelne mellem strømningen i et vandrør og blodstrømmen i årerne, lige så lidt tør vi på forhånd vente nogen dybere principiel forskel mellem sanseindtrykkenes forplantning i nerverne og elektricitetsledningen i en metaltråd. Vel gælder det for alle sådanne fænomener, at en i enkeltheder gående redegørelse fører os ind på atomteoriens område, ja for elektricitetsledningens vedkommende har man netop i de sidste år erkendt, at først den for kvanteteorien ejendommelige begrænsning af de anskuelige bevægelsesforestillinger tillader os at begribe, at elektronerne kan slippe frem mellem metalatomerne i tråden. For disse fænomener gælder det imidlertid, at en sådan uddybelse af beskrivelsen ikke er nødvendig for redegørelsen for de nærmest i betragtning kommende virkninger. Hvad de dybere biologiske problemer angår, hvor det drejer sig om organismens frihed og tilpasningsevne i dens reaktion over for ydre påvirkning, må vi imidlertid regne med, at erkendelsen af en større sammenhæng kan kræve en hensyntagen til de samme forhold, der betinger årsagsbeskrivelsens begrænsning ved atomfænomenerne. Iøvrigt bør vel allerede den kendsgerning, at bevidsthed, som vi kender den, er uadskillelig knyttet til liv, forberede os på, at selve grundproblemet om grænsen mellem levende og dødt kan unddrage sig en forståelse i dette ords sædvanlige forstand. Som undskyldning for, at en fysiker berører sådanne emner, tør måske gælde, at den i fysikken foreliggende, nye situation på så eftertrykkelig måde minder os om den gamle sandhed, at vi såvel er tilskuere som deltagere i tilværelsens store skuespil.
\end{quote}
\dcite{Filosofiske Skrifter 1, Atomteorien og grundprincipperne for naturbeskrivelsen, ¶14 — pp. 82–83 — Bohr's own Danish, matched word for word}
{\footnotesize\begin{quote}\itshape (No corresponding word in the aligned paragraph of the original.)\end{quote}}

\section{Filosofiske Skrifter II}

\subsection{Indledning}

\begin{quote}
DEN fysiske videnskabs betydning for almindelig filosofisk tænkning ligger ikke alene i dens bidrag til vort stadigt voksende kendskab til den natur hvoraf vi selv er del, men tillige i, at den gang på gang har givet anledning til prøvelse og forfinelse af vore begrebsmæssige hjælpemidler. I vort århundrede har studiet af stoffernes atomare opbygning åbenbaret en uanet begrænsning af rækkevidden af den klassiske fysiks forestillinger og derved stillet de i traditionel filosofi overtagne krav til videnskabelig forklaring i nyt lys. Den undersøgelse af forudsætningerne for \tlg{entydig} anvendelse af vore elementære begreber, som sammenfatningen af de atomare fænomener har nødvendiggjort, har derfor bud langt ud over den fysiske videnskabs egentlige område.
\end{quote}
\dcite{Filosofiske Skrifter 2, Indledning, ¶1 — no printed source indexed}

\begin{quote}
Den sidste gruppe af artikler er nært knyttet til den første, men forhåbentlig vil den til fremstillingen af situationen i kvantefysikken benyttede mere træffende terminologi lette forståelsen af den erkendelsesteoretiske indstilling. Ved behandlingen af de mere almindelige problemer er der især lagt eftertryk på forudsætningerne for \tlg{entydig} anvendelse af de ved beskrivelsen af erfaringerne benyttede begreber. Hovedpunktet i argumentationen er fremhævelsen af, at det for objektiv beskrivelse og harmonisk sammenfatning er nødvendigt på næsten ethvert kundskabsområde at være opmærksom på iagttagelsessituationen.
\end{quote}
\dcite{Filosofiske Skrifter 2, Indledning, ¶5 — no printed source indexed}

\subsection{Lys og liv \normalfont(Nature 131 (1933), pp. 421–423 and 457–459)}

\begin{quote}
I de allersidste år har en overordentlig udvikling af atomteorien fundet sted, og vi besidder nu så fuldkomne metoder til at beregne såvel energiværdierne for atomernes stationære tilstande som overgangsprocessernes sandsynlighed, at redegørelsen for atomernes egenskaber efter de af korrespondensargumentet anviste linier med hensyn til omfang og indre sammenhæng næppe står tilbage for den newtonske mekaniks forklaring af astronomiske erfaringer. Selv om den rationelle behandling af atommekanikkens problemer kun har været mulig ved indførelsen af nye symbolske hjælpemidler, er dog den almindelige belæring, som analysen af lysfænomenerne har givet os, stadig bestemmende for vort syn på den omhandlede udvikling. Således består der et komplementaritetsforhold mellem enhver \tlg{entydig} brug af begrebet stationær tilstand og en mekanisk analyse af partiklernes bevægelser inden for atomet, som ganske svarer til det beskrevne forhold mellem lyskvanter og den elektromagnetiske strålingsteori. Ethvert forsøg på at forfølge forløbet af en overgangsproces i enkeltheder ville nemlig medføre en ukontrollerbar energiomsætning mellem atomet og måleinstrumenterne, der fuldstændig ville forstyrre selve den energioverførelse som det var opgaven at gøre rede for. En årsagsbeskrivelse i klassisk forstand kan kun gennemføres i sådanne tilfælde, hvor den virkning det drejer sig om er stor i sammenligning med virkningskvantet, og hvor derfor fænomenerne kan underdeles uden at forstyrres væsentligt. Dersom denne betingelse ikke er opfyldt, kan vi ikke se bort fra vekselvirkningen mellem måleinstrumenterne og undersøgelsesobjektet og må navnlig tage hensyn til, at de forskellige målinger, der ville udkræves for en fuldstændig mekanisk beskrivelse af fænomenerne, kun kan vindes ved forsøgsanordninger, der indbyrdes udelukker hverandre. For helt at forstå denne principielle begrænsning af atomfænomenernes mekaniske analyse må man endvidere gøre sig klart, at det ved en fysisk måling aldrig er muligt at tage vekselvirkningen mellem objektet og måleinstrumenterne direkte i betragtning; i deres egenskaber af iagttagelsesmidler kan instrumenterne nemlig ikke selv samtidig inddrages under undersøgelsen. Ligesom begrebet almindelig relativitet tager sigte på de fysiske fænomeners væsentlige afhængighed af det henførelsessystem, der benyttes til deres indordning i rum og tid, således tjener begrebet komplementaritet til at symbolisere den i atomfysikken afslørede fundamentale begrænsning af vor tilvante forestilling om fænomenerne som eksisterende uafhængigt af de midler hvormed de iagttages.
\end{quote}
\dcite{Filosofiske Skrifter 2, Lys og liv, ¶5 — ≈ p. 422 — aligned to the printed text, not read off a page}
{\footnotesize\begin{quote}\itshape (No corresponding word in the aligned paragraph of the original.)\end{quote}}

\begin{quote}
Ved en nærmere forfølgelse af denne analogi må vi imidlertid ikke glemme, at problemerne frembyder væsentligt forskellige sider i fysikken og i biologien. Medens vi i atomfysikken først og fremmest er interesseret i egenskaberne hos materien i dens simpleste former, er nemlig den indviklede beskaffenhed af de materielle systemer, hvormed vi i biologien har at gøre, af afgørende betydning, da allerede de mest primitive organismer indeholder et stort antal atomer. Vel beror den klassiske mekaniks store anvendelsesområde, som indbefatter vor redegørelse for de måleinstrumenter, der benyttes i atomfysikken, netop på muligheden af ved beskrivelsen af legemer, der indeholder mange atomer, i vid udstrækning at se bort fra den af virkningskvantet betingede komplementaritet. Det er imidlertid typisk for biologiske undersøgelser, at de ydre betingelser, som et særskilt atom er underkastet, aldrig kan kontrolleres på samme måde som i atomfysikkens grundlæggende eksperimenter. Vi kan faktisk ikke engang afgøre, hvilke atomer der strengt taget hører med til en levende organisme, da jo enhver livsfunktion er ledsaget af et stofskifte, ved hvilket atomer stadig optages i og udskilles fra den organisation som udgør det levende væsen. Denne fundamentale forskel mellem fysiske og biologiske undersøgelser fører med sig, at der for de fysiske forestillingers anvendelighed på livsfænomenerne ikke kan trækkes en veldefineret grænse svarende til den inden for atommekanikken gennemførte adskillelse mellem den mekaniske årsagsbeskrivelses område og de egentlige kvanteprocesser. Dog afhænger den begrænsning af den omhandlede analogi, som dette forhold kunne synes at rumme, væsentligt af den måde hvorpå vi finder det bekvemt at bruge sådanne ord som fysik og mekanik. På den ene side ville spørgsmålet om fysikkens begrænsning inden for biologien jo miste enhver mening, dersom vi, i overensstemmelse med ordets oprindelige betydning, ved fysik ville forstå enhver beskrivelse af naturfænomenerne. På den anden side ville et sådant udtryk som atommekanik være vildledende, hvis vi som i daglig sprogbrug kun ville benytte ordet mekanik til at betegne en \tlg{entydig} årsagsbeskrivelse af fænomenerne. Jeg skal ikke her gå nærmere ind på disse rent dialektiske spørgsmål, men endnu kun minde om, at det er det typiske komplementaritetsforhold mellem den underdeling, som en fysisk analyse ville kræve, og sådanne karakteristisk biologiske fænomener som selvopholdelsen og forplantningen af individer, der er kernen i den omhandlede analogi. Det er jo også dette forhold, der medfører at begrebet formålstjenlighed, som ingen plads har i den mekaniske analyse, finder et vist anvendelsesområde ved problemer hvor man må tage hensyn til livets væsen. I denne henseende minder den rolle, som teleologiske argumenter spiller i biologien, om de i korrespondensargumentet formulerede bestræbelser for i atomfysikken at tage virkningskvantet i betragtning på rationel måde.
\end{quote}
\dcite{Filosofiske Skrifter 2, Lys og liv, ¶9 — ≈ p. 458 — aligned to the printed text, not read off a page}
{\footnotesize\begin{quote}\itshape (No corresponding word in the aligned paragraph of the original.)\end{quote}}

\subsection{Biologi og atomfysik \normalfont(Celebrazione del secondo centenario della nascita di Luigi Galvani, Bologna, 18–21 ottobre 1937: I. Rendiconto generale (Tipografia Luigi Parma, 1938), pp. 68–78)}

\begin{quote}
Ethvert spørgsmål om at vende tilbage til en med kausalitetsprincippet overensstemmende beskrivelsesmåde var imidlertid ikke blot udelukket ved det stadigt voksende kendskab til typiske kvanteeffekter; men det viste sig endog snart muligt at udvikle de oprindelige primitive forsøg på i atomteorien at tage virkningskvantets eksistens i betragtning til en principielt statistisk atommekanik, som, med hensyn til konsekvens og fuldstændighed, fuldt tåler sammenligning med den klassiske mekanik, af hvilken den fremtræder som en naturlig almindeliggørelse. Udviklingen af denne nye såkaldte kvantemekanik, som vi som velkendt i første linie skylder en række glimrende bidrag af den yngre generation af fysikere, har, rent bortset fra dens overordentlige frugtbarhed på alle atomfysikkens og kemiens områder, også på væsentlig måde klarlagt det erkendelsesteoretiske grundlag for atomfænomenernes analyse og syntese. Den revision af selve iagttagelsesproblemet på dette område, der indledtes med Heisenbergs opstilling af ubestemthedsprincippet, har faktisk afsløret hidtil upåagtede forudsætninger for den \tlg{entydige} brug af selv de mest elementære begreber, på hvilke naturfænomenernes beskrivelse hviler. Det afgørende punkt er her erkendelsen af, at ethvert forsøg på ved hjælp af den klassiske fysiks metoder og begreber nærmere at analysere den af virkningskvantet betingede »individualitet« i atomprocesserne vil strande på umuligheden af skarpt at skelne mellem vedkommende atomare objekters selvstændige opførsel og deres vekselvirkning med de for en sådan analyse nødvendige måleredskaber.
\end{quote}
\dcite{Filosofiske Skrifter 2, Biologi og atomfysik, ¶12 — ≈ pp. 74–75 — aligned to the printed text, not read off a page; that page is itself inferred}
{\footnotesize\begin{quote}\itshape
Any question of a return to a mode of description consistent with the principle of causality was, however, not only excluded by \olg{unambiguous} experience of the most varied kind, but it soon proved possible to develop the original primitive attempts of taking account of the existence of the quantum of action in atomic theory into a proper, essentially statistical atom mechanics, fully comparable in consistency and completeness with the structure of classical mechanics of which it appears as a rational generakation. … The revision of the very problem of observation in this field, initiated by HEISENBERG, one of the principal founders of quantum mechanics, has in fact led to the disclosure of hitherto disregarded presuppositions for the \olg{unambiguous} use of even the most elementary concepts on which the description of natural phenomena rests.
\end{quote}}

\begin{quote}
Uden på nogen måde at underkende det store opmuntrende forbillede til forfølgelse af sådanne erkendelsesteoretiske undersøgelser, som vi skylder relativitetsteorien, der netop ved at afsløre uanede forudsætninger for alle fysiske begrebers \tlg{entydige} brug åbnede nye muligheder for en harmonisk sammenfatning af tilsyneladende uforenelige fænomener, må vi gøre os klart at den situation, som atomteoriens seneste udvikling har stillet os over for, er helt uden sidestykke i den fysiske videnskabs historie. Hele den klassiske fysiks begrebsbygning, som ved Einsteins værk skulle opnå en så vidunderlig enhed og fuldstændighed, hviler jo på den af vore sædvanlige erfaringer om fysiske fænomener så velbegrundede forudsætning, at det er muligt at skelne mellem de materielle legemers opførsel og vor iagttagelse deraf. For at finde en virkelig parallel til den belæring om sådanne tilvante idealisationers begrænsede gyldighed som atomteorien har givet os, må vi vende os til et fra fysikken så fjernt område af videnskaben som psykologien, eller endog til den art af erkendelsesteoretiske problemer som allerede tænkere som Buddha og Lao Tse stilledes overfor under deres bestræbelser på at finde udtryk for harmonien i tilværelsens store drama, i hvilket vi samtidig er skuespillerne og tilskuerne. Erkendelsen af en sådan lighed i den rent begrebsmæssige karakter af de problemer, som vi møder på så vidt forskellige områder af menneskelig forskning, må dog på ingen måde forveksles med en godkendelse på atomfysikkens område af nogen mod videnskabens sande ånd stridende mystik; tværtimod opmuntrer denne erkendelse os til at undersøge, hvorvidt den logiske løsning af de uforudsete paradokser, som anvendelsen af vore elementæreste begreber på atomfænomenerne frembyder, skulle kunne hjælpe os til også på andre erfaringsområder at overvinde begrebsmæssige vanskeligheder.
\end{quote}
\dcite{Filosofiske Skrifter 2, Biologi og atomfysik, ¶14 — ≈ pp. 75–76 — aligned to the printed text, not read off a page; that page is itself inferred}
{\footnotesize\begin{quote}\itshape (No corresponding word in the aligned paragraph of the original.)\end{quote}}

\begin{quote}
Et sådant synspunkt vil ses at være lige fjernt fra de som mekanisme og vitalisme betegnede yderliggående standpunkter. På den ene side må jo enhver sammenligning mellem levende organismer og maskiner – hvad enten det drejer sig om de gamle »iatro-fysikeres« forholdsvis simple konstruktioner eller de mest forfinede moderne forstærkeranordninger, hvis ukritiske fremhævelse ville gøre os fortjent til øgenavnet »iatro-kvantikere« – fra vort synspunkt forkastes som misvisende. På den anden side må vi tage lige så bestemt afstand fra alle sådanne forsøg på at indføre nogen slags mod velbegrundede fysiske og kemiske lovmæssigheder stridende særegne biologiske love, som vi i vore dage har set genoplivede under det overvældende indtryk af de vidunderlige embryologiske opdagelser vedrørende cellernes vækst og deling. Netop i denne forbindelse må man særlig huske på, at muligheden for inden for komplementaritetssynspunktets ramme at undgå alle sådanne uoverensstemmelser er givet ved selve den kendsgerning, at intet resultat af biologiske undersøgelser kan \tlg{entydigt} beskrives uden ved hjælp af fysiske og kemiske begreber, lige så lidt som vi ved redegørelsen for atomfysikkens erfaringer kan undvære brugen af de for al bevidst indordning af sanseindtryk nødvendige udtryksmidler.
\end{quote}
\dcite{Filosofiske Skrifter 2, Biologi og atomfysik, ¶17 — ≈ pp. 77–78 — aligned to the printed text, not read off a page}
{\footnotesize\begin{quote}\itshape
On the other hand, it rejects as irrational all such attempts at introducing some kind of special biological laws inconsistent with well-established physical and chemical regularities, as have in our days been revived under the impression of the wonderful revelations of embryology regarding cell growth and division, In this connection it must be especially remembered that the possibility of avoiding any such inconsistency within the frame of complementarity is given by the very fact that no result of biological investigation can be \olg{unambiguously} described otherwise than in terms of physics and chemistry, just as any account of experience even in atomic physics must ultimately rest on the use of the concepts indispensable for a conscious recording of sense impressions.
\end{quote}}

\subsection{Fysikkens erkendelseslære og menneskekulturerne \normalfont(Congrès international des sciences anthropologiques et ethnologiques, compte rendu de la deuxième session, Copenhague 1938 (Ejnar Munksgaard, Copenhagen 1939), pp. 86–95)}

\begin{quote}
Hvor gennemgribende en ændring af vort hele syn på naturbeskrivelsen denne udvikling af atomfysikken har ført med sig, belyses måske bedst derved, at endog de rammer, der er afstukne af årsagssætningen som hidtil var betragtet som det uomtvistelige grundlag for enhver forklaring af naturfænomenerne, har vist sig for snævre til at give plads for de egenartede lovmæssigheder, som behersker de individuelle atomprocesser. Jeg behøver næppe at fremhæve, at det har været meget tvingende grunde der har fået fysikerne til at opgive selve kausalitetsidealet; men ved studiet af de nye atomfænomener er vi gang på gang blevet belært om at spørgsmål, som man mente forlængst var blevet endeligt besvarede, rummede uanede overraskelser. De har sikkert alle hørt om de gåder, som de elementæreste problemer vedrørende lysets og materiens egenskaber i de sidste år har stillet os overfor. De tilsyneladende modsigelser, vi i denne henseende har mødt, er i virkeligheden lige så skarpe som de, der i begyndelsen af dette århundrede gav anledning til relativitetsteoriens udvikling, og har ligesom disse først fundet deres opklaring gennem en nærmere undersøgelse af den begrænsning, som de nye erfaringer selv sætter for den \tlg{entydige} anvendelse af de begreber der indgår i fænomenernes beskrivelse. Medens for relativitetsteoriens vedkommende det afgørende punkt var forståelsen af, at to iagttagere, der bevæger sig i forhold til hinanden, vil beskrive givne objekters opførsel på væsentlig forskellig måde, drejer det sig ved opklaringen af atomfysikkens paradokser om erkendelsen af, at den med enhver iagttagelse forbundne uundgåelige vekselvirkning mellem objekterne og måleinstrumenterne endog sætter en principiel grænse for muligheden af overhovedet at kunne tale om en af iagttagelsesmidlerne uafhængig opførsel af atomare objekter.
\end{quote}
\dcite{Filosofiske Skrifter 2, Fysikkens erkendelseslære og menneskekulturerne, ¶4 — ≈ pp. 88–89 — aligned to the printed text, not read off a page}
{\footnotesize\begin{quote}\itshape
The apparent contradictions which we have met in this respect are, in fact, as acute as those which gave rise to the development of the theory of relativity in the beginning of this century and have, just as the latter, only found their explanation by a closer examination of the limitation imposed by the new experiences themselves on the \olg{unambiguous} use of the concepts entering into the description of the phenomena.
\end{quote}}

\subsection{Diskussion med einstein om erkendelsesteoretiske problemer i atomfysikken \normalfont(Albert Einstein: Philosopher-Scientist, ed. P.A. Schilpp, The Library of Living Philosophers, Vol. VII (Evanston, Illinois 1949), pp. 201–241)}

\begin{quote}
Lige fra begyndelsen var diskussionernes hovedpunkt spørgsmålet om, hvordan man skal stille sig til den afvigelse fra tilvante principper for naturbeskrivelsen, som kendetegner den nye udvikling af fysikken, der blev indledt i det første år af dette århundrede med Plancks opdagelse af det universelle virkningskvantum. Denne opdagelse, som afslørede et helhedstræk hos naturlovene der går langt ud over den gamle lære om stoffets begrænsede delelighed, belærte os om at de klassiske fysiske teorier er idealisationer, som kun finder \tlg{entydig} anvendelse i den grænse, hvor alle optrædende virkninger er store i forhold til kvantet. Det omdiskuterede problem har været, hvorvidt det afkald på en årsagsbeskrivelse af atomare processer, som bestræbelserne på at finde sig til rette i situationen havde medført, betød en midlertidig afvigelse fra idealer som til syvende og sidst måtte genindføres, eller om vi stod overfor et uigenkaldeligt skridt til opnåelse af den rette harmoni imellem analyse og syntese af fysiske fænomener. For at beskrive baggrunden for vore diskussioner og så klart som muligt belyse argumenterne for de divergerende synspunkter, har jeg følt det nødvendigt i en vis udstrækning at minde om nogle hovedtræk af udviklingen, til hvilken Einstein selv har ydet så afgørende bidrag.
\end{quote}
\dcite{Filosofiske Skrifter 2, Diskussion med einstein om erkendelsesteoretiske problemer i atomfysikken, ¶2 — ≈ pp. 201–202 — aligned to the printed text, not read off a page}
{\footnotesize\begin{quote}\itshape
This discovery, which revealed a feature of atomicity in the laws of nature going far beyond the old doctrine of the limited divisibility of matter, has indeed taught us that the classical theories of physics are idealizations which can be \olg{unambiguously} applied only in the limit where all actions involved are large compared with the quantum.
\end{quote}}

\begin{quote}
Som det blev betonet i foredraget, frembyder den kvantemekaniske formalisme netop et for den komplementære beskrivelsesmåde egnet matematisk redskab. Vi har jo her at gøre med et rent symbolsk skema, som på korrespondensmæssigt grundlag tillader forudsigelser vedrørende forsøgsresultater der kan opnås under betingelser fastlagt ved hjælp af klassiske begreber. Det må i denne forbindelse erindres, at det selv i ubestemthedsrelationen (3) drejer sig om en konsekvens af formalismen, som det ikke er muligt på \tlg{entydig} måde at fortolke i ord egnede til at beskrive klassiske fysiske billeder. Således giver en sætning som »vi kan ikke kende både bevægelsesmængden og positionen af et atomart objekt« straks anledning til spørgsmål om den fysiske realitet af sådanne attributter til objektet, hvortil svar kun kan gives ved henvisning til betingelserne for den \tlg{entydige} anvendelse af rum-tidsbegreberne på den ene side og de dynamiske bevarelsessætninger på den anden. Medens foreningen af disse begreber i et enkelt billede er grundlaget for den klassiske mekaniks årsagsbeskrivelse, skabes der netop plads til lovmæssigheder uden for denne beskrivelses rækkevidde ved den omstændighed, at studiet af komplementære fænomener kræver forsøgsanordninger, der udelukker hinanden.
\end{quote}
\dcite{Filosofiske Skrifter 2, Diskussion med einstein om erkendelsesteoretiske problemer i atomfysikken, ¶25 — ≈ p. 211 — aligned to the printed text, not read off a page; that page is itself inferred}
{\footnotesize\begin{quote}\itshape
21 I by means of classical concepts, I t must here be remembered that even in the indeterminacy relation (3) we are dealing with an implication of the formalism which defies \olg{unambiguous} expression in words suited to describe classical physical pictures. … Thus, a sentence like “we cannot know both the momentum and the position of an atomic object” raises at once questions as to the physical reality of two such attributes of the object, which can be answered only by referring to the conditions for the \olg{unambiguous} use of space-time concepts, on the one hand, and dynamical conservation laws, on the other hand.
\end{quote}}

\begin{quote}
Nødvendigheden af i atomfysikken påny at undersøge grundlaget for den \tlg{entydige} brug af elementære fysiske forestillinger minder i mange henseender om den situation, som oprindelig førte Einstein til at fremhæve iagttagelsesproblemets afgørende betydning og gav anledning til hans geniale revision af grundlaget for rum-tidskoordinationen, der skulle give vort verdensbillede så stor en enhed. Uagtet alle nye træk ved betragtningsmåden var det muligt i relativitetsteorien at opretholde årsagsbeskrivelsen i ethvert henførelsessystem; i kvanteteorien tvinges vi imidlertid af den ukontrollerbare vekselvirkning mellem objekter og måleinstrumenter til et afkald selv i sådan henseende. Denne erkendelse rummer dog ingenlunde en begrænsning af den kvantemekaniske beskrivelses rækkevidde, og den i Como-foredraget fremførte tankegang tog netop sigte på at vise, at komplementaritetssynspunktet må betragtes som en rationel almindeliggørelse af selve kausalitetsidealet.
\end{quote}
\dcite{Filosofiske Skrifter 2, Diskussion med einstein om erkendelsesteoretiske problemer i atomfysikken, ¶26 — ≈ p. 211 — aligned to the printed text, not read off a page; that page is itself inferred}
{\footnotesize\begin{quote}\itshape
T h e necessity, in atomic physics, of a renewed examination of the foundation for the \olg{unambiguous} use of elementary physical ideas recalls in some way the situation that led Einstein to his original revision on the basis of all application of space-time concepts which, by its emphasis on the primordial importance of the observational problem, has lent such unity to our world picture.
\end{quote}}

\begin{quote}
Det er netop argumenter af den art, der minder os om umuligheden af at underdele kvantefænomener og vanskeligheden ved at tilskrive atomare objekter sædvanlige fysiske attributter. Især må vi gøre os klart, at enhver \tlg{entydig} brug af rum-tidsbegreber i beskrivelsen af atomare fænomener – ud over redegørelsen for placeringen og synkroniseringen af de instrumenter der udgør forsøgsanordningen – er begrænset til registrering af iagttagelser beroende på mærker på en fotografisk plade eller lignende praktisk irreversible forstærkningseffekter, såsom opbygningen af en vanddråbe omkring en ion i et tågekammer. Skønt virkningskvantet selvfølgelig til syvende og sidst er ansvarligt for egenskaberne hos de stoffer, hvoraf måleinstrumenterne er opbygget og på hvilke iagttagelsesmidlernes virkemåde beror, er denne omstændighed ikke relevant for spørgsmålet om hensigtsmæssigheden og fuldstændigheden af de her diskuterede aspekter af den kvantemekaniske beskrivelse.
\end{quote}
\dcite{Filosofiske Skrifter 2, Diskussion med einstein om erkendelsesteoretiske problemer i atomfysikken, ¶52 — ≈ pp. 222–223 — aligned to the printed text, not read off a page}
{\footnotesize\begin{quote}\itshape
I n particular, it must be realized that-besides in the account of the placing and timing of the instruments forming the experimental arrangement-all \olg{unambiguous} use of space-time concepts in the description of atomic phenomena is confined to the recording of observations which refer to marks on a photographic plate or to similar practically irreversible amplification effects like the building of a water drop around an ion in a cloud-chamber.
\end{quote}}

\begin{quote}
Fra vort synspunkt ser vi nu, at formuleringen af det ovenfor nævnte af Einstein, Podolsky og Rosen foreslåede kriterium på fysisk virkelighed rummer en flertydighed med hensyn til meningen af udtrykket »uden på nogen måde at forstyrre et system«. Selvfølgelig er der i et tilfælde som det vi lige har betragtet ikke tale om en mekanisk forstyrrelse af det undersøgte system på det sidste kritiske stadium af målingerne. Men netop på dette stadium er der tale om en indflydelse på selve de betingelser, der definerer de mulige typer af forudsigelser vedrørende systemets fremtidige opførsel. Eftersom disse betingelser udgør et uundværligt element af beskrivelsen af ethvert fænomen til hvilket udtrykket »fysisk virkelighed« på konsekvent måde kan knyttes, ser vi at de nævnte forfatteres argumentation ikke retfærdiggør deres konklusion at den kvantemekaniske beskrivelse er væsentlig ufuldstændig. Denne beskrivelse må tværtimod, som det fremgår af den foregående diskussion, karakteriseres som en rationel udnyttelse af alle muligheder for \tlg{entydig} interpretation afmålinger, forenelig med den af virkningskvantet betingede endelige og ukontrollable vekselvirkning mellem objekterne og måleinstrumenterne på kvanteteoriens område. Det er kun det gensidige udelukkelsesforhold mellem to forsøgsanordninger, tilladende \tlg{entydig} definition af komplementære størrelser, der giver plads for nye fysiske love, hvis optræden ved første blik synes uforenelig med naturvidenskabens grundlæggende principper. Det er netop denne helt nye situation med hensyn til beskrivelse af fysiske fænomener, som betegnelsen komplementaritet tager sigte på at karakterisere.
\end{quote}
\dcite{Filosofiske Skrifter 2, Diskussion med einstein om erkendelsesteoretiske problemer i atomfysikken, ¶79 — ≈ p. 234 — aligned to the printed text, not read off a page}
{\footnotesize\begin{quote}\itshape
O n the contrary, this description, as appears from the preceding discussion, may be characterized as a rational utilization of all possibilities of \olg{unambiguous} interpretation of measurements, compatible with the finite and uncontrollable interaction between the objects and the measuring instruments in the field of quantum theory. … I n fact, it is only the mutual exclusion of any two experimental procedures, permitting the \olg{unambiguous} definition of complementary physical quantities, which provides room for new physical laws, the coexistence of which might at first sight appear irreconcilable with the basic principles of science.
\end{quote}}

\begin{quote}
Omend en sådan indstilling i sig selv kunne synes ganske nøgtern, rummer den ikke desto mindre en afvisning af hele den i det foregående udviklede argumentation, der tager sigte på at vise at vi i kvantemekanikken ikke har at gøre med noget vilkårligt afkald på en mere detaljeret analyse af atomare fænomener, men med erkendelsen af at en sådan analyse er principielt udelukket. Kvanteeffekternes ejendommelige individualitet stiller os med hensyn til sammenfatningen af veldefinerede erfaringer over for en ny situation, der var uforudset i den klassiske fysik og er uforenelig med tilvante forestillinger egnede til vor orientering i og tilpasning til dagliglivets verden. Det er i denne henseende at kvanteteorien har krævet en fornyet revision af grundlaget for den \tlg{entydige} brug af elementære begreber som et videre skridt i den udvikling, der siden relativitetsteoriens fremkomst er blevet så karakteristisk for vor tids videnskab.
\end{quote}
\dcite{Filosofiske Skrifter 2, Diskussion med einstein om erkendelsesteoretiske problemer i atomfysikken, ¶82 — ≈ p. 235 — aligned to the printed text, not read off a page}
{\footnotesize\begin{quote}\itshape
It is in this respect that quantum theory has called for a renewed revision of the foundation for the \olg{unambiguous} use of elementary concepts, as a further step in the development which, since the advent of relativity theory, has been so characteristic of modern science.
\end{quote}}

\begin{quote}
Som en mere træffende udtryksmåde foreslog jeg, at ordet fænomen udelukkende anvendes til at henvise til iagttagelser der er vundet under angivne omstændigheder omfattende en redegørelse for hele forsøgsanordningen. Med en sådan terminologi er iagttagelsesproblemet befriet for en hvilken som helst forvikling, eftersom alle iagttagelser ved faktiske eksperimenter udtrykkes ved \tlg{entydige} udsagn, der f. eks. angår registreringen af det punkt hvor en elektron ankommer til en fotografisk plade. En sådan udtryksmåde er endvidere egnet til at understrege, at det ved den konsekvente fysiske tydning af den symbolske kvantemekaniske formalisme alene drejer sig om forudsigelser, af \tlg{entydig} eller statistisk karakter, vedrørende individuelle fænomener optrædende under betingelser defineret ved klassiske fysiske begreber.
\end{quote}
\dcite{Filosofiske Skrifter 2, Diskussion med einstein om erkendelsesteoretiske problemer i atomfysikken, ¶88 — ≈ pp. 237–238 — aligned to the printed text, not read off a page}
{\footnotesize\begin{quote}\itshape
I n such terminology, the observational problem is free of any special intricacy since, in actual experiments, all observations are expressed by \olg{unambiguous} statements referring, for instance, to the registration of the point at which an electron arrives at a photographic plate.
\end{quote}}

\subsection{Kundskabens enhed}

\begin{quote}
FØR vi forsøger at besvare spørgsmålet om, i hvilken udstrækning man kan tale om kundskabens enhed, må vi spørge om betydningen af selve ordet kundskab. Det er dog ikke min hensigt at holde et akademisk filosofisk foredrag, hvortil jeg næppe besidder den fornødne lærdom. Enhver videnskabsmand bliver imidlertid til stadighed stillet over for problemet om objektiv beskrivelse af erfaringerne, hvormed vi simpelthen mener \tlg{entydige} meddelelser. Vort fundamentale redskab er selvfølgelig det daglige sprog, der tjener det praktiske livs og det sociale samkvems tarv, og vi skal her ikke beskæftige os med dette sprogs oprindelse, men med dets rækkevidde ved videnskabelig meddelelse og især med problemet om, hvordan beskrivelsens objektivitet kan bevares, når erfaringsområdet overskrider dagliglivets hændelser.
\end{quote}
\dcite{Filosofiske Skrifter 2, Kundskabens enhed, ¶1 — no printed source indexed}

\begin{quote}
Når vi taler om en begrebsmæssig ramme, henviser vi blot til utvetydig logisk fremstilling af relationer mellem erfaringer. Denne indstilling kommer også til udtryk i den historiske udvikling, idet man ikke længere skelner skarpt mellem formel logik og studier af semantik eller endog filologisk syntaks. En særlig rolle spiller selvfølgelig matematikken, som har bidraget så afgørende til udviklingen af logisk tænkning og hvis veldefinerede abstraktioner er så uvurderlig en hjælp til at udtrykke harmonisk sammenhæng. I vor diskussion skal vi dog ikke betragte matematikken som en særskilt kundskabsgren, men snarere som en forfinelse af det fælles sprog, der forsyner dette med hjælpemidler egnet til at fremstille relationer, for hvis meddelelse sædvanlige sproglige udtryk enten ikke er tilstrækkelig skarpe eller er for ubekvemme. I denne forbindelse må det understreges, at netop ved at undgå sådan henvisning til det bevidste subjekt, som i så høj grad gennemsyrer dagligsproget, tjener brugen af matematiske symboler til at sikre den af objektiv beskrivelse krævede \tlg{entydighed} af definitionerne.
\end{quote}
\dcite{Filosofiske Skrifter 2, Kundskabens enhed, ¶3 — no printed source indexed}

\begin{quote}
Udviklingen af de såkaldte eksakte videnskaber, der er kendetegnet ved opstillingen af numeriske sammenhæng mellem målinger, er blevet afgørende fremmet ved benyttelsen af abstrakte matematiske metoder, der ofte var udformet uden henblik på sådan anvendelse, men alene skyldtes bestræbelsen på at almindeliggøre logiske konstruktioner. Denne situation er især illustreret i fysikken, hvorved man oprindelig forstod al kundskab om den natur, hvoraf vi selv er del, men som efterhånden kom til at betyde studiet af de elementære love for egenskaberne hos det livløse stof. Nødvendigheden af selv inden for dette forholdsvis simple emne til stadighed at være opmærksom på problemet om objektiv beskrivelse har gennem tiderne haft dyb indflydelse på de filosofiske skolers indstilling. I vore dage har udforskningen af nye erfaringsområder afsløret upåagtede forudsætninger for \tlg{entydig} anvendelse af vore mest elementære begreber og derved givet os en erkendelsesteoretisk belæring af betydning for problemer langt uden for fysikkens område. For vor diskussion vil det derfor være bekvemt at begynde med en kort redegørelse for denne udvikling.
\end{quote}
\dcite{Filosofiske Skrifter 2, Kundskabens enhed, ¶4 — no printed source indexed}

\begin{quote}
Inden for sit store anvendelsesområde udgør den klassiske mekanik en objektiv beskrivelse i den forstand, at den hviler på veldefineret brug af billeder og forestillinger, der viser hen til dagliglivets hændelser. Hvor rationelle de i den newtonske mekanik benyttede idealisationer end syntes, gik de dog langt ud over rækkevidden af de erfaringer, til hvilke vore elementære begreber er tilpasset. Hensigtsmæssigheden af selve begreberne absolut rum og tid er således uløseligt forbundet med lysets praktisk talt øjeblikkelige forplantning, der tillader at lokalisere legemerne omkring os uafhængigt af deres hastigheder og at ordne begivenhederne i en \tlg{entydig} tidsrække. Bestræbelsen på at udvikle en modsigelsesfri redegørelse for elektromagnetiske og optiske fænomener åbenbarede imidlertid, at iagttagere der bevæger sig i forhold til hverandre med store hastigheder vil indordne begivenhederne på forskellig måde. Sådanne iagttagere vil ikke alene bedømme faste legemers form og størrelse forskelligt, men begivenheder i forskellige punkter i rummet, som for én iagttager fremtræder som samtidige, vil af en anden bedømmes som følgende efter hverandre.
\end{quote}
\dcite{Filosofiske Skrifter 2, Kundskabens enhed, ¶6 — no printed source indexed}

\subsection{Atomerne og den menneskelige erkendelse}

\begin{quote}
Ved en hvilken som helst forsøgsanordning, der tillader at konstatere en partikels tilstedeværelse i et begrænset rum-tidsområde, må der jo benyttes fast anbragte målestokke og synkroniserede ure, der efter deres definition udelukker kontrol af en til dem overført bevægelsesmængde eller energi. Omvendt kræver enhver \tlg{entydig} anvendelse af dynamiske bevarelseslove i kvantefysikken, at det omhandlede fænomens beskrivelse for de atomare objekters vedkommende indebærer et principielt afkald med hensyn til detaljeret rum-tidskoordination. Dette gensidige udelukkelsesforhold mellem de forsøgsbetingelser, som de elementære begrebers brug forudsætter, medfører at hele forsøgsanordningen må tages i betragtning ved fænomenernes veldefinerede beskrivelse. I denne sammenhæng finder kvantefænomenernes udelelighed konsekvent udtryk deri, at enhver definerbar underdeling ville kræve en ændring af forsøgsanordningen med optræden af nye individuelle fænomener. Selve grundlaget for en deterministisk beskrivelse er således faldet bort, og forudsigelsernes statistiske karakter fremtræder derved, at der under samme forsøgsomstændigheder i almindelighed fremkommer iagttagelser svarende til forskellige mulige individuelle processer.
\end{quote}
\dcite{Filosofiske Skrifter 2, Atomerne og den menneskelige erkendelse, ¶17 — no printed source indexed}

\begin{quote}
På baggrund af den indflydelse, som den mekaniske naturopfattelse har udøvet på filosofisk tænkning, er det forståeligt at man fra mange sider har opfattet komplementaritetssynspunktet som rummende en med beskrivelsens objektivitet uforenelig henvisning til den subjektive iagttager. Selvfølgelig må vi på ethvert erfaringsområde opretholde en skarp adskillelse mellem iagttageren og indholdet af iagttagelserne, men vi må betænke at virkningskvantets opdagelse har stillet selve grundlaget for naturbeskrivelsen i ny belysning og belært os om hidtil upåagtede forudsætninger for den rationelle anvendelse af de begreber, på hvilke meddelelserne om erfaringerne hviler. I kvantefysikken er, som vi har set, en redegørelse for måleinstrumenternes funktioner uundværlig for definitionen af fænomenerne, og vi må så at sige drage skillelinien mellem subjekt og objekt på en måde, der i hvert enkelt tilfælde sikrer den \tlg{entydige} anvendelse af de i meddelelserne benyttede elementære fysiske begreber. Langtfra at rumme en mod videnskabens ånd stridende mystik henviser betegnelsen komplementaritet blot til de med vor stilling ved beskrivelsen og sammenfatningen af erfaringerne på atomfysikkens område forbundne erkendelsesvilkår.
\end{quote}
\dcite{Filosofiske Skrifter 2, Atomerne og den menneskelige erkendelse, ¶19 — no printed source indexed}

\subsection{Fysikken og livets problem \normalfont(Atomic Physics and Human Knowledge (John Wiley \& Sons, New York 1958), pp. 94–101)}

\begin{quote}
For at karakterisere forholdet mellem fænomener iagttagne under forskellige forsøgsomstændigheder har man indført betegnelsen komplementaritet for at minde om, at sådanne fænomener til sammen udtømmer al definerbar oplysning om de omhandlede atomare objekter. Langt fra at rumme noget vilkårligt afkald på tilvant fysisk forklaring henviser komplementaritetssynspunktet direkte til vor stilling som iagttagere på et erfaringsområde, hvor den \tlg{entydige} anvendelse af de begreber, der benyttes ved fænomenernes beskrivelse, væsentlig afhænger af iagttagelsesomstændighederne. Ved en matematisk almindeliggørelse af den klassiske fysiks begrebsbygning er det lykkedes at udvikle en formalisme, hvori virkningskvantet indgår på harmonisk måde. Denne såkaldte kvantemekanik tager direkte sigte på formuleringen af statistiske lovmæssigheder gældende for erfaringer vundne under veldefinerede iagttagelsesbetingelser, og beskrivelsens principielle fuldstændighed beror på opretholdelsen af de klassiske mekaniske forestillinger i en udstrækning, der omfatter enhver definerbar variation af forsøgsomstændighederne.
\end{quote}
\dcite{Filosofiske Skrifter 2, Fysikken og livets problem, ¶15 — ≈ p. 99 — aligned to the printed text, not read off a page}
{\footnotesize\begin{quote}\itshape
Far from containing any arbitrary renunciation of customary physical explanation, the notion of complementarity refers directly to our position as observers in a domain of experience where \olg{unambiguous} application of the concepts used in the description of phenomena depends essentially on the conditions of observation.
\end{quote}}

\begin{quote}
Sammenlignet med den udvidelse af den mekanistiske beskrivelsesmåde, som redegørelsen for de atomare fænomeners helhedstræk har krævet, stiller organismens integritet og personlighedens enhed os over for en langt mere vidtgående almindeliggørelse af rammerne for vore meddelelsesmidlers rationelle brug. I denne henseende må det fremhæves, at den for \tlg{entydig} beskrivelse nødvendige skelnen mellem subjekt og objekt er opretholdt derved, at vi ved enhver meddelelse som rummer en henvisning til os selv så at sige indskyder et nyt subjekt, der ikke optræder som led i meddelelsens indhold. Det behøver næppe at understreges, at det netop er denne frihed i valget af subjekt-objekt skillelinjen, der giver plads for bevidsthedsfænomenernes mangfoldighed og menneskelivets muligheder.
\end{quote}
\dcite{Filosofiske Skrifter 2, Fysikken og livets problem, ¶21 — ≈ p. 101 — aligned to the printed text, not read off a page; that page is itself inferred}
{\footnotesize\begin{quote}\itshape
In this respect, it must be emphasized that the distinction between subject and object, necessary for \olg{unambiguous} description, is retained in the way that in every communication containing a reference to ourselves we, so-to-speak, introduce a new subject which does not appear as part of the content of the communication.
\end{quote}}

\part*{IV.\quad Filosofiske Skrifter III--IV (counts only)}
\addcontentsline{toc}{part}{IV.\quad Filosofiske Skrifter III--IV}

\noindent These volumes translate Bohr's English and German; \emph{entydig} there is the translators' word, not his. Counted, not quoted.

\begin{center}\begin{tabular}{@{}p{0.8\textwidth}r@{}}\toprule essay & paragraphs\\\midrule
FS 3: Kvantefysik og filosofi & 1\\
FS 3: Den menneskelige erkendelses enhed & 2\\
FS 3: Rutherford mindeforelæsning 1958 & 3\\
FS 3: Kvantemekanikkens tilblivelse & 2\\
FS 3: Solvay-møderne og kvantefysikkens udvikling & 2\\
FS 4: Faraday lecture: kemi og kvanteteorien om atomernes opbygning & 11\\
FS 4: Kvantemekanik og fysisk virkelighed & 1\\
FS 4: Kan den kvantemekaniske beskrivelse af fysisk virkelighed opfattes som fuldstændig? & 7\\
FS 4: Kausalitet og komplementaritet & 5\\
FS 4: Analyse og syntese i videnskaben & 1\\
FS 4: Kausalitetsproblemet i atomfysikken & 9\\
FS 4: Om begreberne kausalitet og komplementaritet & 4\\
FS 4: Fysisk videnskab og studiet af religioner & 2\\
FS 4: Fysisk videnskab og menneskets stilling & 4\\
\bottomrule\end{tabular}\end{center}

\end{document}
